\section{Data transformation and knowledge-graph population}
\label{sec:transform}
The transformation is implemented in Python. Apache Jena Fuseki is not the transformer: it
is the RDF store and query server used after the graph has been created. Figure~\ref{fig:transform-flow}
shows the implemented path from source records to the validated serving union.

\begin{figure}[htbp]\centering
\resizebox{\textwidth}{!}{\begin{tikzpicture}[
  node distance=8mm and 7mm,
  tstep/.style={draw=accent!80!black, rounded corners=3pt, fill=accent!5, align=center, font=\small, text width=3.15cm, minimum height=14mm},
  toutput/.style={draw=blue!70!black, rounded corners=3pt, fill=blue!7, align=center, font=\small, text width=3.15cm, minimum height=14mm},
  arr/.style={-{Stealth}, thick, draw=accent!80!black},
  aux/.style={-{Stealth}, dashed, thick, draw=black!55}
]
\node[tstep] (s1) at (0,0) {1. Collect\\API responses, infoboxes,\\cached source revisions};
\node[tstep] (s2) at (4.0,0) {2. Integrate\\normalize labels, resolve\\identifiers and conflicts};
\node[tstep] (s3) at (8.0,0) {3. Transform\\mint persistent IRIs,\\emit typed RDF};
\node[tstep] (s4) at (12.0,0) {4. Enrich\\external links, OWL rules,\\qualified observations};
\node[toutput] (s5) at (16.0,0) {5. Validate and release\\SHACL, profile, metrics,\\fingerprints and serving union};

\node[toutput] (a1) at (2.0,-2.55) {Bronze JSON\\HTTP cache};
\node[toutput] (a2) at (6.0,-2.55) {Silver JSON\\conflict and unresolved logs};
\node[toutput] (a3) at (10.0,-2.55) {vnedu-data.ttl\\URI registry};
\node[toutput] (a4) at (14.0,-2.55) {vnedu-links.ttl\\vnedu-inferred.ttl};
\node[toutput] (a5) at (16.0,-4.75) {vnedu-all.ttl\\HTML / Turtle / JSON-LD\\Fuseki input};

\draw[arr] (s1)--(s2)--(s3)--(s4)--(s5);
\draw[aux] (s1)--(a1);
\draw[aux] (s2)--(a2);
\draw[aux] (s3)--(a3);
\draw[aux] (s4)--(a4);
\draw[arr] (s5)--(a5);
\node[font=\scriptsize, align=center] at (8,-1.0) {each step is reproducible from versioned inputs; a failed gate blocks publication};
\end{tikzpicture}
}
\caption[Data transformation workflow]{Data transformation workflow. Each stage produces a
reviewable artifact. The final serving union is emitted only after RDF, ontology, reasoning,
SHACL, metric, and fingerprint checks succeed.}
\label{fig:transform-flow}
\end{figure}

\subsection{Bronze layer: collection and source preservation}
The Bronze layer stores cached responses, structured snapshots, and source context. Wikidata
queries contribute identifiers and structured attributes; Vietnamese Wikipedia contributes
infobox fields, article text, links, and revision identifiers; reference tables contribute
administrative crosswalks, fields, majors, and illustrative programs. A cache miss fails by
default during reconstruction, so a normal build does not silently change its input by
contacting a live endpoint. Fresh collection is an explicit operation.

\subsection{Silver layer: integration and normalization}
The integration stage joins records through stable identifiers when available, normalizes
labels and values, resolves organization references, and records conflicts or unresolved
targets. It produces normalized JSON for institutions, people, governing bodies, provinces,
and programs. The Silver layer is where source policies are applied; it is not yet an OWL
graph. For example, a ministry associated with a private institution is represented as
reported state oversight rather than automatically asserted as direct ownership or governance.

\subsection{Gold layer: RDF transformation, linking and reasoning}
The RDF transformer maps each normalized record to a class and a set of properties in
\texttt{ontology/vnedu.ttl}. Persistent IRIs are reused from \texttt{uri\_map.json}; a readable
slug is only the initial allocation key. Literal values receive explicit language tags and
datatypes, while resource-valued fields become object IRIs. External links are written to a
separate link graph. The reasoner applies the local ontology and selected OWL/RDF rules, and
the release retains selected instance-level entailments rather than treating every internal
closure statement as an authoritative source fact.

\subsection{Concrete Silver-to-RDF mapping}
Table~\ref{tab:mapping} follows the shipped Silver record \texttt{Q3075696}, Hanoi University
of Science and Technology. It demonstrates the mapping rules; it is not an independent
verification of the institution's current legal status or leadership.

\begin{table}[htbp]\centering\small
\caption{Selected Silver-to-RDF mappings for one institution record.}\label{tab:mapping}
\begin{tabularx}{\textwidth}{p{3.3cm}Xp{4.4cm}}\toprule
Silver input & RDF output & Rule or qualification \\\midrule
\texttt{key = Q3075696} & Subject \texttt{u:dai-hoc-bach-khoa-ha-noi} & Reuse the persistent registry; a label change does not mint a new IRI. \\
\texttt{kind = University} & \texttt{rdf:type vnedu:University} & Classification follows the integration policy, not Fuseki. \\
\texttt{name\_vi} & \texttt{rdfs:label} and \texttt{skos:prefLabel} with \texttt{@vi} & Preserve language information. \\
\texttt{founding\_year = 1956} & \texttt{vnedu:foundingYear ``1956'' \textasciicircum\textasciicircum xsd:integer} & May refer to predecessor origin; the event meaning remains qualified. \\
\texttt{province = Q1858} & \texttt{vnedu:locatedIn} the registered Hanoi IRI & Use an object resource, not a province-name literal. \\
\texttt{ownership = public} & \texttt{vnedu:ownership vnedu:PublicOwnership} & Supports the public-institution classification rule. \\
\texttt{viwiki\_revid} & \texttt{prov:wasDerivedFrom} a Wikipedia oldid IRI & Entity-level source context, not complete assertion-level provenance. \\
\bottomrule
\end{tabularx}
\end{table}

Missing optional values normally produce no triple; they are not converted into empty strings,
zeros, or claims of non-existence. Reference programs are identified using the institution slug
and major code. This gives a stable demonstration resource, but it does not distinguish every
curriculum, degree level, delivery mode, or academic year without additional identifiers.

\subsection{Release boundary}
The serving graph is the set union of ontology, asserted data, selected inference, external
links, and metadata. Component files remain available for inspection, but loading the union
into a default graph does not recreate named-graph provenance. The workflow is therefore a
reproducible staged batch process, not a transaction covering every intermediate file. The
validated release record binds the accepted output to its inputs and rejects stale publication
or Fuseki loading attempts.
